\section{Results}
\label{sec:results}

The results below come from bench trials run by the author on the
physical dVRK, used to validate that the pipeline works as described
in Section~\ref{sec:methods}, not from a study with trainees; that
study is future work (Section~\ref{sec:conclusion}). Rather than
validating each measure in isolation, this chapter follows the same
logic as the pipeline itself (Section~\ref{subsec:pipeline_overview}):
the measures are merged into one system, so results are shown as what
one trial, run through the whole pipeline at once, actually produces,
with timing broken out on its own, since it calls for a direct,
quantitative check against the real-time claim this thesis makes.

\subsection{An Integrated Trial}
\label{subsec:results_integrated_trial}

% PLACEHOLDER: walk through one representative trial end to end, in
% the order the pipeline itself produces it (Section~\ref{subsec:pipeline_overview}):
% trial time and path length live in the GUI; rate of orientation
% change printed to the terminal; the upright/color-down check and
% object-state classification at placement; the force popup after the
% trial ends. The point is to show these together, as the trainee
% would actually see them from a single trial, not as separate
% one-off numbers.
% TODO: pick 1-3 representative trials (e.g. a clean success, a drop,
% a wrong-color placement) and report what every measure showed for
% each, ideally with a screenshot/photo of the GUI and force popup.

\subsection{Timing}
\label{subsec:results_timing}

This thesis claims the pipeline is real-time. This section checks
that claim directly, per module, against measured timing rather than
against the code's own assumptions about itself.

Trial time, path length, and rate of orientation change
(Section~\ref{subsec:kinematic_metrics}) update continuously while the
clock runs, driven by a background process that assumes a 200~Hz
internal clock (Section~\ref{subsec:gui_trial_clock}). The measured
update rate, both how often the GUI actually refreshes and how often
new PSM pose samples actually arrive, was [TODO] Hz.

The orientation check (Section~\ref{subsec:angular_displacement}) has
a built-in 0.5~s settle delay after the Arduino's ``Target hit''
signal, plus processing time, plus up to two further 0.5~s retries if
the cylinder is not found in the crop. Measured from that event to the
printed result, over [TODO] trials, the time was [TODO]~s on average
and [TODO]~s at the slowest.

The object-state classifier (Section~\ref{subsec:object_state}) runs
once per frame; its measured inference time, over [TODO] calls, was
[TODO]~ms on average and [TODO]~ms at the slowest.

The force popup (Section~\ref{subsec:ground_contact}) appears once a
trial ends; the measured time from trial end to the popup becoming
visible was [TODO]~s.
